\section{Variables de Decisión: Estabilidad y Eficiencia}
\label{sec:variables}

El análisis utiliza dos propiedades moleculares y una medida final de proceso.
La primera indica cuánto resiste la enzima; la segunda, qué tan bien convierte;
la tercera, si esa mejora realmente llega a producción.

\subsection{Variable X: Termoestabilidad}
\label{subsec:termoestabilidad}

La termoestabilidad mide cuánto tiempo una enzima conserva actividad bajo una
temperatura determinada. En este trabajo se observa mediante dos indicadores:

\begin{enumerate}
    \item \(t_{50}\): tiempo necesario para que la actividad disminuya al 50\%;
    \item área bajo la curva de actividad residual: actividad conservada durante
    todo el periodo evaluado.
\end{enumerate}

Un valor mayor en ambos indicadores implica una vida operativa más extensa.
Esta propiedad es crítica cuando la enzima se inactiva antes de completar la
conversión \citep{siddiqui_defying_2017,hou_tradeoff_2023}.

\subsection{Variable Y: Eficiencia Catalítica}
\label{subsec:eficiencia}

La eficiencia catalítica indica qué tan eficazmente la enzima transforma el
sustrato disponible. Se estima a partir de la relación entre velocidad de
reacción y concentración de sustrato.

Para la comparación se utiliza la relación:

\begin{equation}
\text{Eficiencia Catalítica}
=
\frac{V_{\max}}{K_m},
\end{equation}

donde \(V_{\max}\) representa la velocidad máxima estimada y \(K_m\) caracteriza
la respuesta de la enzima frente al sustrato.

Además, se considera que una concentración excesiva puede reducir la velocidad.
Por ello, más sustrato no se interpreta automáticamente como una mejora del
proceso.

\subsection{Variable Final: Producción}
\label{subsec:produccion_variable}

La medida industrial final es el producto acumulado durante ocho horas. Esta
variable permite comprobar si una mejora en estabilidad o eficiencia se
convierte efectivamente en mayor producción.

\subsection{Mapa de Variables}
\label{subsec:mapa}

La Figura \ref{fig:variables} resume la lógica utilizada.

\begin{figure}[H]
\centering

\resizebox{0.98\textwidth}{!}{
\begin{tikzpicture}[
    node distance=1.0cm and 1.0cm,
    caja/.style={
        draw,
        rounded corners=5pt,
        minimum width=3.7cm,
        minimum height=1.25cm,
        align=center,
        inner sep=6pt
    },
    flecha/.style={-{Latex[length=3mm]},thick}
]

\node[caja] (x) {Termoestabilidad\\\(t_{50}\) y actividad residual};
\node[caja, below=1.0cm of x] (y) {Eficiencia catalítica\\velocidad frente a sustrato};

\node[caja, right=1.4cm of x, yshift=-1.05cm] (proceso)
{Condiciones de operación\\temperatura, sustrato, producto};

\node[caja, right=1.4cm of proceso] (resultado)
{Resultado industrial\\producto acumulado};

\draw[flecha] (x.east) -- (proceso.west);
\draw[flecha] (y.east) -- (proceso.west);
\draw[flecha] (proceso) -- (resultado);

\end{tikzpicture}
}

\figcaptionlink{anx:fig_variables}{Relación entre las Variables del Análisis}{fig:variables}


\end{figure}

\subsection{Criterio Central}
\label{subsec:criterio_xy}

Existe trade-off cuando una mejora de estabilidad viene acompañada de una caída
de eficiencia. Sin embargo, esta relación no es obligatoria
\citep{siddiqui_defying_2017,hou_tradeoff_2023}.

La condición de mayor interés es:

\begin{equation}
\text{Mayor Estabilidad}
\quad+\quad
\text{Mayor Eficiencia}
\quad\longrightarrow\quad
\text{Mayor Producción}.
\end{equation}

La \hyperref[sec:metodologia]{Sección 3} define cómo se cuantifica cada mejora.
